% Auto-drafted from runs/all_metrics_summary.csv on the current experiment branch.
% These paragraphs cover completed E2--E5 results only.  E1 cross-model
% diagnostics should be added after GPT-4o/Gemini/Qwen FinEvo runs export.

\subsection{Component Ablations}
\label{sec:component_ablations}

Table~\ref{tab:component_ablation} reports the five-seed component ablation in
the 10-agent, 24-month setting.  The full FinEvo configuration reaches
$182.1\mathrm{k}\pm10.7\mathrm{k}$ final wealth with low unemployment
($0.17\%\pm0.37\%$) and no invalid or failed API calls.  Removing episodic
memory causes the clearest labor-market degradation, increasing unemployment to
$6.67\%\pm2.57\%$ and reducing final wealth to
$116.8\mathrm{k}\pm9.5\mathrm{k}$.  Removing reflection also substantially
reduces wealth ($104.4\mathrm{k}\pm15.1\mathrm{k}$) and increases inequality
($0.428\pm0.083$ Gini), indicating that the reflection-to-rule pathway changes
subsequent decisions rather than serving only as an interpretability layer.
Removing the sentiment gate produces a smaller but consistent wealth reduction,
while removing semantic memory leaves wealth similar to the full model but
raises inequality.  We therefore treat the modules as contributing along
different axes: episodic memory is most important for avoiding repeated labor
failures, reflection is most important for wealth and inequality, and semantic
rules primarily affect distributional stability.

\subsection{Sentiment Parameterization Robustness}
\label{sec:sentiment_robustness}

Table~\ref{tab:sentiment_robustness} tests whether FinEvo depends on a single
hand-specified sentiment formula.  We perturb each coefficient by $\pm50\%$,
remove macroeconomic terms, and replace the sentiment channel with a random
control.  Across coefficient perturbations and the no-macro control, final
wealth remains in a relatively narrow range
($169.7\mathrm{k}$--$191.0\mathrm{k}$), unemployment remains at or below
$0.33\%$, and all runs have zero invalid-action and API-error rates.  The random
sentiment control is the main exception, reducing wealth to
$145.3\mathrm{k}\pm23.2\mathrm{k}$.  These results suggest that the exact
default coefficients are not the only load-bearing factor, while a semantically
unstructured sentiment signal can still degrade outcomes.

\subsection{Textual Event Channel Controls}
\label{sec:text_event_controls}

Table~\ref{tab:text_event_controls} separates numeric-only sentiment from the
synthetic textual event channel.  The aligned text-event condition performs
similarly to numeric-only sentiment, while the shuffled-event control reduces
wealth to $136.4\mathrm{k}\pm10.7\mathrm{k}$.  The random-event condition does
not dominate the aligned condition and introduces a small unemployment rate
($0.33\%\pm0.46\%$).  These controls address the ambiguity in the earlier
implementation, where the news shock was a numerical disturbance rather than an
explicit text event.  The current evidence should be interpreted as a synthetic
event-fixture control, not as a claim about real-time news understanding.

\subsection{Prompt and Decoding Robustness}
\label{sec:prompt_robustness}

Table~\ref{tab:prompt_robustness} evaluates whether the observed behavior is
driven by prompt wording, decoding temperature, or output format.  The concise,
risk-neutral, no-fairness, and temperature variants remain in the same
qualitative range as the default prompt, and all variants have zero invalid
actions and zero API errors.  The natural-language-plus-JSON condition produces
lower wealth ($162.2\mathrm{k}\pm28.5\mathrm{k}$), indicating that looser output
formatting can affect behavior even when the parser succeeds.  We therefore
report output compliance together with economic metrics and avoid attributing
cross-model differences solely to economic preferences.
